\section{Follow-up studies G2 to G4: details}\label{app:I}

This appendix holds the tables behind Sections~\ref{sec:stop} and~\ref{sec:ctc} and the B8 warm-start paragraph of Section~\ref{sec:lm-oracle}.
All three studies are oracle-memory studies at rank 128 on the 768 seen names with one oracle seed.
Their values are produced by \texttt{scripts/paper\_tables\_g2.py} from the committed files of each study (\texttt{reports/lsr}, \texttt{reports/b8w}, \texttt{reports/ctc}), which checks its key values against the figures the three study reports state.
Rounding is half up on the stored value, so a cell that ends in an exact 5 (for example 64.55) prints 0.1 higher than in the study reports.
The evidence ledger holds the exact values (rows LS1 to LS4, BW1 to BW3, CT1 to CT3) and wins over any shortened wording here.

\subsection{LSR: label-free stop rules (G2)}

The test range is [106000, 108000), 2,000 episodes at eight hops per checkpoint, and the six checkpoints are the registered rank-128 M1-L4 runs at Qwen3-0.6B-Base (P1-L) and Qwen3-1.7B-Base (P5L), seeds 17, 29, 43.
Thresholds were chosen on validation episodes [1000, 3000) and nothing was retrained.
The registered intervals are 99.17\% paired episode-bootstrap intervals (six comparisons, 10,000 draws shared across seeds, hops and rules); descriptive intervals are 95\%.  % LS1
Table~\ref{tab:lsr} in Section~\ref{sec:stop} gives the EM by hop of the fixed-depth rows, R2 and the two depth-per-hop references.
Table~\ref{tab:i-lsr} gives R2 against fixed L=4.

\begin{table}[htbp]
\centering\footnotesize
\caption{R2 against fixed L=4 (which equals the best fixed depth L\textsubscript{c} under the validation constraint), pp.  % LS1 LS3
Intervals are 99.17\% for the registered rows (hops 4 to 8, hop 1) and 95\% for the descriptive rows.}
\label{tab:i-lsr}
\begin{tabular}{@{}W{400}{4}W{300}{4}W{300}{4}@{}}
\toprule
 & Qwen3-0.6B-Base & Qwen3-1.7B-Base \\
\midrule
Hops 4 to 8 gain & +3.90 [3.35, 4.43] & +5.68 [5.09, 6.25] \\ \addlinespace[3pt]
Hops 4 to 8 gain per seed (17/29/43) & +4.58/+3.25/+3.87 & +5.46/+5.71/+5.86 \\ \addlinespace[3pt]
Hop 1 change & +6.28 [5.09, 7.51] & +4.70 [3.48, 5.92] \\ \addlinespace[3pt]
Hop 1 change per seed & +8.90/+3.95/+6.00 & +7.35/+4.40/+2.35 \\ \addlinespace[3pt]
Hop 2 change (descriptive) & $-10.83$ [$-11.85$, $-9.87$] & $-12.27$ [$-13.32$, $-11.23$] \\ \addlinespace[3pt]
Hop 3 change & $-3.28$ & $-3.65$ \\ \addlinespace[3pt]
Hops 1 to 3 pooled change & $-2.61$ [$-3.18$, $-2.06$] & $-3.74$ [$-4.32$, $-3.16$] \\ \addlinespace[3pt]
R2 minus L=8, hops 4 to 8 & $-3.54$ [$-3.93$, $-3.15$] & $-5.05$ [$-5.47$, $-4.62$] \\ \addlinespace[3pt]
R2 minus L=8, hop 1 & +34.1 & +39.4 \\ \addlinespace[3pt]
R2 minus L=4 by hop 1 to 8 & +6.28, $-10.83$, $-3.28$, +4.48, +4.98, +4.27, +3.08, +2.68 & +4.70, $-12.27$, $-3.65$, +7.35, +7.85, +5.32, +4.33, +3.53 \\ \addlinespace[3pt]
Mean reads per question & 5.60 & 5.66 \\ \addlinespace[3pt]
Upper-layer passes per question ($s^{*}+1$) & 6.60 & 6.66 \\ \addlinespace[3pt]
Mean reads by hop 1 to 8 & 4.97, 5.48, 5.67, 5.76, 5.67, 5.70, 5.73, 5.86 & 4.87, 5.46, 5.72, 5.87, 5.80, 5.77, 5.84, 5.93 \\
\bottomrule
\end{tabular}
\end{table}

The registered conditions (gain of at least 3 pp with an interval above 0, hop 1 not below $-10$ pp, every seed) are all met at both backbones.  % LS1
R1 and R3 are infeasible at both backbones: none of their 25 validation grid values kept the hop-1 loss against fixed L=4 within 8 pp pooled and 10 pp per seed, so their test statistics were never computed.  % LS2
A prediction recorded before registration (R3: hops 4 to 8 at least 12\%, hop 1 at least 55\% at 0.6B) was not a gate and is not met.  % LS2
Latency is descriptive and was measured within one validation job, in ms per question (0.6B/1.7B): fixed L=4 0.56/0.62, fixed L=8 0.74/0.81, the L=8 trace pass 0.81/0.87.
R2 itself was not timed, because its early exit needs a pass through the upper layers for each lens evaluation, and no latency or efficiency statement is made.  % LS3
At hops 4 to 8 R2 (9.1\% and 10.2\%) lies between fixed L=4 (5.2\% and 4.5\%) and fixed L=8 (12.6\% and 15.3\%), and the validation-chosen depth per hop (O1) reaches 13.0\% and 15.8\%.  % LS3 LS4
A fixed L=5 reaches 8.5\% on hops 4 to 8 at 0.6B with hop 1 at 36.6\%, close to R2's 9.1\% with hop 1 at 54.3\%; it lies outside the validation constraint (it lost more than the hop-1 margin) and no interval is attached to this comparison.  % LS4
The intervals condition on three training seeds per backbone and one oracle seed and layout, and they do not cover new training seeds.

\subsection{B8W: a B8 copy started from M1-L4 (G3)}

The continuations start from the registered P1-L M1-L4 checkpoints (Qwen3-0.6B-Base, rank 128, seeds 17, 29, 43) and run 2,000 steps at learning rate 0.003 with a fresh optimizer on the P1-L stream, in two arms: the shared M1-L4 continued, and a B8 copy with four untied step matrices.
Evaluation is on test [108000, 110000), 2,000 episodes, L=4 (B4 at L=1), with 95\% paired episode-bootstrap intervals (2,000 draws).
The equality gate (bf16 logits, predicted tokens, the control and the fp32 read) held in every seed with a maximum absolute difference of 0 at tolerance 0, and the protocol check (shared arm more than 2 pp below the start) was not triggered.  % BW1
Table~\ref{tab:i-b8w} gives every registered reference contrast, and Table~\ref{tab:i-b8w-hop} the accuracy by hop.

\begin{table}[htbp]
\centering\footnotesize
\caption{B8W contrasts on hops 2 to 3 unless stated (pp, 95\% interval).  % BW1 BW2 BW3
The first row is the registered decision; the others are descriptive.}
\label{tab:i-b8w}
\begin{tabular}{@{}W{330}{4}W{250}{4}W{190}{4}W{230}{4}@{}}
\toprule
Contrast & Difference [95\% interval] & Accuracies (\%) & Per seed (17/29/43) \\
\midrule
untied@4 $-$ shared@4 (decision) & +0.90 [0.09, 1.71] & 32.40 vs 31.50 & +1.05/+0.90/+0.75 \\ \addlinespace[3pt]
untied $-$ shared, hop 1 & +2.77 [1.55, 4.02] & 55.47 vs 52.70 & +2.85/+1.70/+3.75 \\ \addlinespace[3pt]
untied $-$ start (M1-L4@4) & +5.99 [5.07, 6.88] & 32.40 vs 26.41 & +5.08/+6.38/+6.53 \\ \addlinespace[3pt]
shared $-$ start & +5.09 [4.18, 6.01] & 31.50 vs 26.41 & +4.03/+5.48/+5.78 \\ \addlinespace[3pt]
untied $-$ start, hop 1 & +7.45 [6.00, 8.95] & 55.47 vs 48.02 & +6.40/+8.75/+7.20 \\ \addlinespace[3pt]
shared $-$ start, hop 1 & +4.68 [3.33, 6.02] & 52.70 vs 48.02 & +3.55/+7.05/+3.45 \\ \addlinespace[3pt]
untied $-$ original B8@4 (P1-L) & +26.07 [24.98, 27.13] & 32.40 vs 6.33 & +29.38/+20.03/+28.80 \\ \addlinespace[3pt]
untied $-$ B4@1 & +26.98 [25.87, 28.08] & 32.40 vs 5.42 & +29.88/+21.38/+29.70 \\ \addlinespace[3pt]
shared $-$ B4@1 & +26.08 [24.98, 27.19] & 31.50 vs 5.42 & +28.83/+20.47/+28.95 \\ \addlinespace[3pt]
start M1-L4@4 $-$ B4@1 & +20.99 [19.97, 22.06] & 26.41 vs 5.42 & +24.80/+15.00/+23.18 \\
\bottomrule
\end{tabular}
\end{table}

The decision rule asks that every seed loses at most 2 pp and that the pooled lower bound is above $-2$ pp; both hold, and no seed has a loss at all.  % BW1
The 97.5\% interval of the decision contrast is [$-0.04$, 1.81], which touches 0.
No superiority test was registered, so the untied arm being nominally above the shared arm is exploratory and is not a finding.  % BW2
Both arms improved by about 5 to 6 pp over the start, so the starting reader had not converged, and this says nothing about untying being better.  % BW2 BW3

\begin{table}[htbp]
\centering\footnotesize
\caption{B8W: exact match (\%) by hop, seed means, L=4 (B4 at L=1).  % BW3
}
\label{tab:i-b8w-hop}
\begin{tabular}{@{}lrrrrrrrr@{}}
\toprule
Reader & 1 & 2 & 3 & 4 & 5 & 6 & 7 & 8 \\
\midrule
Start (M1-L4) & 48.0 & 34.7 & 18.2 & 6.8 & 4.6 & 5.1 & 4.9 & 5.1 \\
Shared, continued & 52.7 & 41.9 & 21.1 & 5.5 & 3.6 & 4.2 & 3.7 & 4.1 \\
Untied copy & 55.5 & 40.8 & 24.0 & 4.2 & 3.4 & 4.0 & 3.6 & 3.6 \\
Original B8 (P1-L) & 18.8 & 6.7 & 6.0 & 5.8 & 5.9 & 5.9 & 6.2 & 6.0 \\
B4@1 & 20.0 & 5.6 & 5.3 & 5.1 & 5.6 & 5.3 & 5.5 & 5.1 \\
\bottomrule
\end{tabular}
\end{table}

The starting M1-L4, the original B8 and B4@1 on hops 2 to 3 (26.41\%, 6.33\%, 5.42\%) match the values of P1-L on its own range (26.91\%, 6.34\%, 5.36\%) within 0.5 pp.
At hops 4 to 8 both continued arms are below the starting reader, which is descriptive.  % BW3
The relative Frobenius change of the continued tensors from their start is, for R, 0.44 to 0.49 (shared) and 0.48 to 0.58 (untied copies); for U, 0.22 to 0.23 and 0.26 to 0.32; for Q, 0.21 to 0.24 and 0.22 to 0.29.
The mean pairwise cosine distance between the four untied copies is 0.191 to 0.215 for R, 0.072 to 0.074 for U and 0.051 for Q.
So the copies separated from each other, and the untied arm did not stay tied in practice.  % BW3
Not claimed: that sharing is or is not necessary for composition, the cause of B8's failure from a random start (optimization, update allocation and the inductive bias of sharing are not separated; the original B8 stays at 6.33\% on hops 2 to 3), a revised G1-L, equal-parameter status for B8, or anything at rank 32, for learned writers, unseen names, other scales or other training lengths.  % BW2
The study used 1.004 node-hours over ten jobs.

\subsection{CTC: a classify-then-code writer (G4)}

All values are on validation episodes [0, 1000) with evaluation seed 4242, one seed, full-vocabulary EM.
The readers are the frozen P1-L rank-128 checkpoints of training seed 17 (B3-L4, B4, M1-L4).
Table~\ref{tab:i-ctc} gives M1-L4 at L=4 for every memory condition.
The hard lookup writes the code of the predicted class; the soft lookup writes the class-probability mixture of the codes; the two differ by at most 0.1 pp at every hop for both arms.  % CT1

\begin{table}[htbp]
\centering\footnotesize
\caption{CTC: M1-L4 at L=4, exact match (\%) by hop on validation episodes [0, 1000).  % CT1 CT3
CE and CE+ans are the two registered training arms (class cross-entropy only; with an answer loss added).}
\label{tab:i-ctc}
\begin{tabular}{@{}lrrrrrrrr@{}}
\toprule
Memory & 1 & 2 & 3 & 4 & 5 & 6 & 7 & 8 \\
\midrule
Oracle memory (reference, $p=1$) & 50.6 & 36.6 & 21.5 & 5.9 & 2.2 & 4.2 & 3.5 & 4.6 \\
CE, hard (selected) & 50.7 & 36.5 & 21.6 & 6.0 & 2.2 & 4.2 & 3.6 & 4.5 \\
CE, soft & 50.6 & 36.6 & 21.4 & 5.9 & 2.2 & 4.3 & 3.7 & 4.7 \\
CE+ans, hard & 50.7 & 36.5 & 21.7 & 6.0 & 2.2 & 4.2 & 3.6 & 4.6 \\
CE+ans, soft & 50.8 & 36.5 & 21.7 & 6.1 & 2.3 & 4.4 & 3.6 & 4.7 \\
Zero memory & 0.0 & 0.0 & 0.0 & 0.0 & 0.1 & 0.0 & 0.0 & 0.1 \\
Shuffled memory (CE) & 0.2 & 0.0 & 0.2 & 0.3 & 0.2 & 0.1 & 0.0 & 0.3 \\
\bottomrule
\end{tabular}
\end{table}

The per-fact joint classification accuracy is 1.00000 for CE (12,000 of 12,000 facts, every episode fully correct) and 0.99992 for CE+ans (one fact wrong), and a perfect-prediction lookup reproduces the oracle memory bit for bit on these episodes.  % CT1
The gain over B4@1 on hops 2 to 3 is +24.05 pp [21.80, 26.30] for CE hard against +23.90 pp with the oracle memory; hop 1 is 50.7\% against a threshold of 40\%.
Neither registered falsifier fires, and the arm choice is a tie (0.0 pp), which the rule gives to CE.  % CT1
B4 at L=1 gives 18.7/5.3/4.7\% on hops 1 to 3 (CE hard) against 18.5/5.4/4.9\% with the oracle memory, and 4.5 to 6.0\% on hops 4 to 8 in both; with zero memory it gives 0.2/0.1/0.2\%.
The memory-free B3-L4 gives 0.0\% at every hop, and its predictions are identical under the oracle and the zero memory.  % CT3
The writer is not a learned writer in the sense of the DPFP rows: both entity heads read their entity at its own token, and the relation is constant in this task, so the relation head scores 100\% at every position.  % CT1 CT2
The study shows that the reader transfers to near-oracle codes and uses its memory; it does not show that a model can learn its own codes, and no CTC number is placed beside a DPFP number.  % CT2
Arm CE trained 8,000 steps in 533 s (0.0651 s per step) and arm CE+ans in 2,320 s (0.2883 s per step), because it also runs the frozen reader and its backward pass every step; the study used 1.028 node-hours against 1.8 registered.  % CT3
